\section{Discussion}\label{sec:discussion}

\subsection{Interpretation and methodological limits}
The strongest methodological lesson is that observing an outcome and receiving its report are different events. Calibration needs the latter, while valid evaluation uses the former. Keeping these channels separate exposes a failure mode that input-only missingness experiments cannot fully characterise: the forecasting system can lose both explanatory observations and the evidence needed to correct its uncertainty.

The results favour a conditional interpretation. The asymmetric redesign succeeds on fresh SUMO C3 and transfers consistently across METR-LA backbones, while PEMS-BAY remains heterogeneous. Residual shape, graph neighbourhoods, scale quality and network dynamics may contribute, but the experiments do not isolate a dominant explanation. The failed symmetric confirmations also caution against attributing all final gains to a single component change: successive formulations were evaluated on different held-out stages.

Full ablations sharpen this conclusion. Stale blending is repeatedly useful, and its removal under permanent loss can collapse coverage. Conversely, target-age and context weighting sometimes hurt score. These findings suggest a future programme of training/tuning-only component selection followed by validation on new networks or untouched episode sets. Choosing a revised rule from the current test ablations would be further method development, not confirmation of the locked rule.

Neither weighted pooling nor the support statistic supplies a missing-outcome identification argument. Selective feedback can change the residual population, especially if low speeds are preferentially suppressed. The MNAR experiment quantifies sensitivity to a specified fault generator; it cannot recover the distribution of unseen outcomes without assumptions or additional information. There is no claim of identifiable MNAR correction or finite-sample coverage under arbitrary reporting failures.

\subsection{Motivation, societal implications and deployment boundary}
The practical aim is to make uncertainty services more accountable when sensing infrastructure degrades. Wider or better-positioned intervals may help operators recognise that a point forecast is less informative during an outage. That possibility is a motivation, not a demonstrated improvement in routing, travel time, safety, emissions or public welfare. No traffic-control policy is evaluated in this study.

Uneven detector availability also raises distributional concerns. Regions with persistently poor instrumentation may receive intervals built from neighbouring or historical residuals that do not represent local conditions. Aggregate network coverage can hide such disparities. A deployment should therefore audit coverage and interval quality by location, reporting quality and relevant operational regime, communicate outages explicitly and retain human oversight for consequential actions. The present datasets and simulation are insufficient to claim equitable performance across communities or road types.

The architecture records forecast issuance and subsequent feedback, which can support auditing. A production implementation would additionally require access controls, retention policies, operational monitoring and validation on live infrastructure. The existing gateway tests establish replay-equivalence properties at smoke scale; they do not replace a full service study. This manuscript accordingly claims a causal replay benchmark and tested calibration behaviour, not operational readiness.

\subsection{Limitations}
The real-data audit is hourly, fault focused and locally capped. Its conclusions cannot be substituted for complete five-minute network-wide evaluation. Simulation provides controlled replication but represents a small freeway network. CoRel, STID and DCRNN are documented adaptations, and STAEformer is an independent implementation; the comparisons do not reproduce every source-paper setting. The archive-cap study uses one fixed run per network and descriptive timing. Ablation intervals are exploratory and not adjusted for multiple comparisons. Dataset redistribution rights for the exact real benchmark packages remain unverified, so the release provides acquisition instructions rather than redistributing them. Finally, the literature positioning is bounded: no universal absence of related feedback-aware work is asserted.


\subsection{What the comparisons establish}
The backbone comparison and the component study answer complementary questions. Table~\ref{tab:backbones} shows that a method can transfer across forecasting representations without performing uniformly across networks. Tables~\ref{tab:ablation-metr}--\ref{tab:ablation-sumo} show that a useful full pipeline can contain components that are adverse in specific cells. Together they argue for conditional rather than universal recommendations. The experiments do not identify whether scale quality, residual asymmetry, neighbourhood heterogeneity or temporal dynamics drives the PEMS-BAY interaction. Establishing that mechanism would require an additional registered comparison, not a post hoc explanation presented as fact.

The fresh C3 simulation validation offers protection against selecting a candidate from its own validation outcomes. C4 transfer is a separate assessment on the existing held-out set and misses the coverage criterion. These distinctions matter because increasing the number of model or fault seeds does not create new independent physical episodes. Likewise, the 432 ablation stages are an execution count, not 432 independent inferential observations. The paired resampling units remain episodes in simulation and day blocks within fault seeds in the real audit.

\subsection{Relationship to recent feedback theory}
The 2026 delayed-feedback and corrupted-feedback preprints reinforce the need to treat outcome reporting as part of uncertainty inference \citep{elhalabi2026,wang2026}. They also restrict the claim of originality. Our contribution is not a first treatment of feedback problems or a stronger coverage theorem. It is the traffic-specific executable intervention design, its archive-based signed-residual kernel and its completed controls, transfer and removal evidence. The published incomplete-time-series study also overlaps with the motivation \citep{chen2026}. Only explicit differences in the reviewed protocols are asserted; absence of a feature in an accessible source is not proof that no related work exists.

\subsection{Implications for subsequent studies}
A next confirmatory study could preregister network-specific eligibility for age and context weighting, evaluate it on newly held-out networks and compare calibration under equal resource budgets. Sensor-level and fault-regime coverage would test whether aggregate performance hides local failures. A separately registered live study would be required for latency and operational reliability claims, while a traffic-control experiment would be required for travel-time or safety benefits. These proposals follow from observed limitations; none is counted as completed evidence in this manuscript.
